\documentclass[12pt,a4paper]{article}
\usepackage[utf8]{vietnam}
\usepackage{amsmath,amssymb}
\usepackage{geometry}
\geometry{
  top=12mm,
  bottom=12mm,
  left=14mm,
  right=14mm,
  footskip=6mm,
  headheight=14pt
}

\usepackage{tikz}
\usetikzlibrary{arrows.meta,calc,shapes.geometric}
\usepackage[most]{tcolorbox}
\usepackage{fontawesome}
\usepackage{fancyhdr}
\usepackage{enumitem}

\pagestyle{fancy}
\fancyhf{}
\fancyhead[L]{\footnotesize\color{teal!80!black}\faBook\ \textbf{SỔ TAY KIẾN THỨC TOÁN 3}}
\fancyhead[R]{\footnotesize\color{teal!80!black}\textbf{Lớp Toán Thầy Lĩnh -- loptoanthaylinh.vn}}
\fancyfoot[C]{\footnotesize\color{gray}-- Sổ tay kiến thức: Bài 15. Luyện tập chung bảng nhân, bảng chia --}
\renewcommand{\headrulewidth}{0.4pt}
\renewcommand{\footrulewidth}{0pt}

\setlength{\parindent}{0pt}
\setlength{\parskip}{0pt}

\begin{document}

% --- BANNER BÀI HỌC ---
\begin{tcolorbox}[
  enhanced,
  colback=teal!8!white,
  colframe=teal!70!black,
  arc=3mm, boxrule=1.2pt,
  left=3mm, right=3mm, top=2mm, bottom=2mm,
  halign=center
]
  {\LARGE\bfseries\color{teal!90!black}\faBookmark\ BÀI 15: LUYỆN TẬP CHUNG BẢNG NHÂN, BẢNG CHIA}\\[1mm]
  {\itshape\footnotesize\color{black!75} Đại tổng kết các bảng nhân chia từ 2 đến 9, tính giá trị biểu thức và giải toán có lời văn trên 1 trang A4}
\end{tcolorbox}

\vspace{-2mm}

\noindent
\begin{minipage}[t]{0.49\textwidth}\vspace{0pt}
\begin{tcolorbox}[
  enhanced,
  colback=blue!4!white,
  colframe=blue!75!black,
  arc=2mm, boxrule=0.8pt,
  left=2mm, right=2mm, top=1.8mm, bottom=1.8mm,
  title={\footnotesize\bfseries\color{white}\faKey\ 1. Ma trận Bảng nhân chia 2 - 9}
]
{\footnotesize
\begin{itemize}\setlength{\itemsep}{1pt}\setlength{\leftskip}{-2.5mm}
  \item \textbf{Mối quan hệ nhân chéo:}
  $$a \times b = c \iff c : a = b \iff c : b = a$$
  \item \textbf{Tính chất đặc biệt với 0 và 1:}
  \begin{itemize}\setlength{\itemsep}{0.5pt}\setlength{\leftskip}{-2.5mm}
    \item $a \times 1 = 1 \times a = a$; $a : 1 = a$.
    \item $a \times 0 = 0 \times a = 0$; $0 : a = 0\ (a \ne 0)$.
    \item Tuyệt đối không có phép chia cho số $0$!
  \end{itemize}
\end{itemize}
}
\end{tcolorbox}
\end{minipage}
\hfill
\begin{minipage}[t]{0.49\textwidth}\vspace{0pt}
\begin{tcolorbox}[
  enhanced,
  colback=teal!4!white,
  colframe=teal!75!black,
  arc=2mm, boxrule=0.8pt,
  left=2mm, right=2mm, top=1.8mm, bottom=1.8mm,
  title={\footnotesize\bfseries\color{white}\faCalculator\ 2. Thứ tự thực hiện phép tính}
]
{\footnotesize
\begin{itemize}\setlength{\itemsep}{1pt}\setlength{\leftskip}{-2.5mm}
  \item \textbf{Biểu thức có cả $(+, -)$ và $(\times, :)$:}
  $$\text{Nhân / Chia TRƯỚC} \implies \text{Cộng / Trừ SAU}$$
  \item \textbf{Biểu thức chỉ có $(\times, :)$:} Thực hiện lần lượt từ \textbf{trái sang phải}.
  \item \textbf{Có dấu ngoặc tròn $( )$:} Tính toán trong ngoặc trước, ngoài ngoặc sau.
\end{itemize}
}
\end{tcolorbox}
\end{minipage}

\vspace{-2.5mm}
% --- KHỐI 3: SƠ ĐỒ TƯ DUY TIKZ ---
\noindent
\begin{tcolorbox}[
  enhanced,
  colback=orange!4!white,
  colframe=orange!80!black,
  arc=2mm, boxrule=0.8pt,
  left=2mm, right=2mm, top=1.5mm, bottom=1.5mm,
  title={\footnotesize\bfseries\color{white}\faSitemap\ SƠ ĐỒ TƯ DUY: HỆ THỐNG 4 MŨI NHỌN TÍNH TOÁN BẢNG NHÂN CHIA}
]
\centering
\begin{tikzpicture}[
  box/.style={draw, rounded corners=2pt, align=center, font=\scriptsize\bfseries, inner sep=2.5pt, minimum height=0.75cm},
  arr/.style={-{Stealth[length=2mm]}, thick}
]
  \node[box, fill=purple!15, draw=purple!80!black, text width=2.6cm] (b1) at (0, 0) {BẢNG NHÂN 2 - 9\\Thuộc lòng các mốc\\Giao hoán $a \times b = b \times a$};
  \node[box, fill=teal!15, draw=teal!80!black, text width=2.6cm] (b2) at (3.5, 0) {BẢNG CHIA 2 - 9\\Nhẩm từ phép nhân\\Tìm 1 phần mấy $\to$ Chia};
  \node[box, fill=blue!12, draw=blue!80!black, text width=2.6cm] (b3) at (7.0, 0) {TÌM THÀNH PHẦN\\$x \times a = b \to x = b : a$\\$a : x = b \to x = a : b$};
  \node[box, fill=orange!15, draw=orange!80!black, text width=2.6cm] (b4) at (10.5, 0) {THỨ TỰ PHÉP TÍNH\\Nhân chia trước\\Cộng trừ sau};

  \draw[arr, purple!80!black] (b1) -- (b2);
  \draw[arr, teal!80!black] (b2) -- (b3);
  \draw[arr, blue!80!black] (b3) -- (b4);
\end{tikzpicture}
\end{tcolorbox}

\vspace{-2.5mm}
% --- KHỐI 4: VÍ DỤ & MẸO TRÁNH BẪY ---
\noindent
\begin{minipage}[t]{0.49\textwidth}\vspace{0pt}
\begin{tcolorbox}[
  enhanced,
  colback=green!4!white,
  colframe=green!60!black,
  arc=2mm, boxrule=0.8pt,
  left=2mm, right=2mm, top=1.8mm, bottom=1.8mm,
  title={\footnotesize\bfseries\color{white}\faCheckCircle\ Ví dụ thực hành chuẩn mực}
]
{\footnotesize
\textbf{1. Tính giá trị biểu thức:}\\
$6 \times 7 + 18 = 42 + 18 = \mathbf{60}$.\\
$72 : 8 - 4 = 9 - 4 = \mathbf{5}$.\\
$8 \times (15 - 9) = 8 \times 6 = \mathbf{48}$.\\
\textbf{2. Tìm $x$:}\\
$x \times 7 = 63 \implies x = 63 : 7 = \mathbf{9}$.\\
$56 : x = 8 \implies x = 56 : 8 = \mathbf{7}$.
}
\end{tcolorbox}
\end{minipage}
\hfill
\begin{minipage}[t]{0.49\textwidth}\vspace{0pt}
\begin{tcolorbox}[
  enhanced,
  colback=red!4!white,
  colframe=red!70!black,
  arc=2mm, boxrule=0.8pt,
  left=2mm, right=2mm, top=1.8mm, bottom=1.8mm,
  title={\footnotesize\bfseries\color{white}\faExclamationTriangle\ Mẹo tránh bẫy trắc nghiệm}
]
{\footnotesize
\begin{itemize}\setlength{\itemsep}{1.5pt}\setlength{\leftskip}{-2.5mm}
  \item \textbf{Bẫy tính từ trái sang phải mà không nhìn phép nhân:}\\
  $10 + 2 \times 5$: Rất nhiều bạn lấy $10 + 2 = 12$ rồi nhân $5$ được $60$ (SAI!). Phải tính $2 \times 5 = 10$ trước, rồi $10 + 10 = \mathbf{20}$.
  \item \textbf{Bẫy nhân chia liên tiếp:}\\
  $18 : 3 \times 2$: Phải tính $18 : 3 = 6$, sau đó $6 \times 2 = \mathbf{12}$. Không được lấy $3 \times 2 = 6$ trước.
\end{itemize}
}
\end{tcolorbox}
\end{minipage}


\end{document}
